\section{Perplexity：最传统也最容易误读的指标}

语言模型定义了 token 序列 \(x\) 的概率 \(p(x)\)。给定数据集 \(D\)，perplexity 可写为：
\[
\mathrm{PPL}(D)
  =
\left(\frac{1}{p(D)}\right)^{1/|D|}
  =
\exp\left(-\frac{1}{|D|}\sum_{x_i\in D}\log p(x_i)\right).
\]
直觉上，PPL 约等于模型在每个 token 位置“等效面对多少个选项”。PPL 越低，模型越不惊讶。但它只测概率分配，不直接测“回答是否有用”。

\begin{importantbox}{术语消化：perplexity 的边界}
Perplexity 适合做平滑 scaling-law 指标，因为它连续、低噪声、可在预训练数据上自然计算。它不适合单独代表真实产品质量，因为它会惩罚所有 token，包括与任务无关的表达形式；也不能直接评价工具使用、开放式帮助、安全性和 agent 行为。
\end{importantbox}

\begin{knowledgebox}{术语澄清：zero-shot 不是 ZeRO}
本节的 zero-shot 指“不针对目标数据集微调，直接评测”。它不是 ZeRO（Zero Redundancy Optimizer）：ZeRO 是训练系统里的优化器状态、梯度和参数分片方法，用于减少 data-parallel ranks 上的重复显存。
\end{knowledgebox}

\begin{figure}[H]
\centering
\includegraphics[width=0.9\textwidth]{gpt2-perplexity.png}
\caption{GPT-2 zero-shot perplexity：WebText 训练后在标准数据集上做 out-of-distribution evaluation。}
\figsource{CS336 Lecture 12 官方源引用图。}
\end{figure}

\begin{knowledgebox}{读图：GPT-2 perplexity 说明 transfer 不是均匀发生}
GPT-2 在小数据集上可能因为预训练 transfer 表现更好，但在更大或分布更不同的数据集上不一定占优。这说明 in-distribution perplexity 和 out-of-distribution perplexity 是不同问题。读这张图时，不要把“某个数据集上 PPL 下降”泛化成“所有语言任务都更好”。它更像在问：WebText 训练出的概率分布，能否迁移到这些标准数据集？
\end{knowledgebox}

老师随后提醒，许多看起来像 task benchmark 的东西，其实仍然是概率比较。LAMBADA 和 HellaSwag 把 perplexity 包装成填空或多选题，让评价更接近任务，但底层仍在比较候选 continuation 的 likelihood。

\begin{figure}[H]
\centering
\includegraphics[width=0.82\textwidth]{lambada.png}
\caption{LAMBADA：cloze task，填补最后一个词，属于 perplexity in disguise。}
\figsource{CS336 Lecture 12 官方源引用图。}
\end{figure}

\begin{figure}[H]
\centering
\includegraphics[width=0.62\textwidth]{hellaswag.png}
\caption{HellaSwag：multiple-choice sentence completion，也可看作比较候选 completion 的概率。}
\figsource{CS336 Lecture 12 官方源引用图。}
\end{figure}

\begin{knowledgebox}{课堂提示：为什么说它们是 perplexity in disguise}
LAMBADA 和 HellaSwag 都可以通过给候选 continuation 打 log probability 来选择答案。它们比纯 PPL 更接近任务，但仍依赖模型概率是否校准、tokenization 是否公平、候选长度如何归一化。老师在这里的隐含提醒是：任务形式变了，不代表评价对象完全变了。
\end{knowledgebox}

\begin{warningbox}{perplexity leaderboard 的信任问题}
如果参赛者提交的是 black-box LM，评测方计算 \(\log p(\text{test data})\) 时必须相信模型返回的是合法概率分布并且归一化。下游任务只需 response，而 perplexity leaderboard 需要概率接口，安全边界不同。
\end{warningbox}

\subsection{本章小结}

Perplexity 仍是语言模型开发中最常用的内部指标，因为它平滑、便宜、适合 scaling。可是它离真实使用很远。接下来考试类 benchmark 试图补足“能否答题”的维度，但又会牺牲开放性和真实感。

